\section{Practical analysis}
\label{sec:practical_analysis}

Since \snowman{} is similar to \bhtb{} and \snortmanager{}, it was
important to benchmark their performance in various formal test cases in order to
highlight the strengths and weaknesses in each software solution. The
purpose of the tests is to reach final conclusions on key questions regarding
the performance of \snowman{}, and to point out where the difference between the
three software packages lies.

To get a picture of how the software performs on different systems, it was
tested on three architectures (table \ref{tab:testSystems}). To minimize error, each test
was run five times where the mean value of the runs is considered a valid test
result. All tests can be reproduced with the test procedure documentation in
Appendix \ref{app:test_scripts}.

\subsection{Definitions}

The testing is split up into \textit{test suites}, each containing a
\textit{problem} in the form of a question. The problem forms a clear
description of what is to be tested, and how results should be evaluated.
A test suite also contains one or more \textit{test cases}, which exist to
address the problem. Test cases represent individual tests, and describe the input data
and constraints. Test cases are carried out on the different systems, and
results from each test case are referred to as the \textit{test results}.

\subsection{Environment}
%Beskriv testoppsett. maskinvare, nettverk, programvare, etc.
All tests are carried out on three systems, named \textit{S1}, \textit{S2} and
\textit{S3} (Table \ref{tab:testSystems}). These systems differ in processing
power, available memory and storage technology, and the results will indicate
what impact these factors have on software performance. The performance
testing in this report is also prone to error, as tests might get disturbed by
other background processes etc. Testing on three systems gives results with a
higher confidence, as results are checked aginst each other for 
inconsistencies. When an inconsistency in results is found, the test is run
again to see if the inconsistency prevails.

%Set row spacing:
\renewcommand{\arraystretch}{1.5}

\begin{table}[h!]
  \centering
  \begin{tabular}{l|l|l|l}
  	%& \multicolumn{1}{c|}{Test} &
  	%\multicolumn{1}{c|}{S1} & \multicolumn{1}{c|}{S2} & \multicolumn{1}{c}{S3}
  	% \\
	System & Processor type & Memory & Storage \\
	\hline
	S1 & Intel\textregistered{} Core\texttrademark{} 2 Solo, 1.4GHz & 4GB RAM & 5K
	RPM HDD \\
	\hline
	S2 & Intel\textregistered{} Core\texttrademark{} i7, 6x3.2GHz & 12GB RAM & 10K
	RPM HDD \\
	\hline
	S3 & Intel\textregistered{} Core\texttrademark{} i5, 4x3.2GHz & 16GB RAM & SSD
	\\
	\hline
  \end{tabular}
  \caption{Test systems.}
  \label{tab:testSystems}
\end{table}

\subsection{Test suites}
Even though \snowman{}, \bhtb{} and \snortmanager{} share the same goal of being a
rule management system, it is not immediately obvious how they can be tested against
each other to compare their performance. The first step in the performance
testing was to identify the areas where the three software packages share, to a
certain extent\footnote{\snowman{}, \bhtb{} and \snortmanager{} have been built
to tackle many of the same challenges, however, the actual implementations may
differ severely. For an extended walk-through of differences in implementation,
refer to Chapter \ref{chap:existing_solutions}, and Chapter
\ref{chap:analysis_of_results}.}, the same functionality. In the
test planning, considerations have been made to ensure that the tests hold an
acceptable degree of fairness (i.e. tests are designed so that none of the
programs enjoys any improper advantages or disatvantages), and that the tests
are pertinent to the daily operation of the software and the \snowman{}
requirement specification.

\snowman{}, \bhtb{} and \snortmanager{} all have the functionality to download and
process an update from an external source. The first group of test
suites presents different test cases to benchmark how each software performs in
update processing.

\subsection{Update processing tests}
The purpose of the tests in this section is first and foremost to test how
quickly the software can process an update, and if the processing time is
sublinear regarding the update size. It is interesting to test for sublinearity
as this may help point out programming errors. To get a complete picture of the
performance, the tests are set up with different constraints such as update
size, number of files and number of changes from one update to the next.

Tests numbered 1-1 through 1-4 use an artificially generated update package
where all rules contain the same data\footnote{Only signature IDs are different. The
rules contain a message and a classtype.}, as seen in Listing
\ref{lst:test_rulestring}.

\begin{lstlisting}[frame=single, caption=Test rule string.,
label=lst:test_rulestring,captionpos=b]
alert tcp $EXTERNAL_NET any -> $HOME_NET $HTTP_PORTS 
(msg:"INDICATOR-COMPROMISE c99shell.php 
command request - security";
flow:to_server,established; urilen:<50; content:"act=security";
fast_pattern:only; http_uri; metadata:service http; 
classtype:policy-violation; sid:<int>; rev:2;)
\end{lstlisting}
%Set row spacing:
\renewcommand{\arraystretch}{1.5}

\begin{table}[h!]
  \centering
  \begin{tabular}{l|l|l}
  	\multicolumn{1}{l|}{Test case} & \multicolumn{1}{|l|}{In test} &
  	\multicolumn{1}{l}{Description}\\
	\hline
		SYN-1 & Test 1-1, 1-2 & Blank database. \\ 
		SYN-10 && \\
		SYN-100&&\\
		SYN-1000&&\\
		SYN-10000&&\\
	\hline
		SYN-1000-B & Test 1-3 & All rules in one file.\\
	\hline		
		SYN-1000-C & Test 1-4 & Pre-populated database. No changes in update. \\
	\hline
		SYN-1000-D & Test 1-5 & Pre-populated database. Revision change for all rules
		\\
	\hline
		SF-28064 & Test 1-1 & Blank database. \\
	\hline
		SF-28064-B & Test 1-6 & Pre-populated database. Update has one rule change.\\
	\hline
	\multicolumn{3}{l}{\textit{Note: Rules are split into files with ten rules per
	file unless otherwise specified.}}
  \end{tabular}
  \caption{Test cases in update processing.}
  \label{tab:updateTests}
\end{table}

\subsubsection{Test 1-1}
\paragraph{Problem: How fast can the software process an update?}
The purpose of this question is to determine the average processing time per
rule. Results are ranked from short to long processing time, where shorter is
better.
\paragraph{Test cases:}
To address this problem, five test cases were introduced. The tests contained an
artificially generated update with 1, 10, 100, 1000 and 10000 rules,
respectively. Rules were split into files each containing ten rules. Each test
case represents one update, and the names denote the number of rules in the following structure:
\textit{SYN-number of rules}. The test was also carried out with a realistic
update package, \textit{SF-28064}.

\subsubsection{Test 1-2}
\paragraph{Problem: How strong is the linear correlation between processing time
and the size of an update?}
This question rises from the requirements specification of Snowman (Chapter
\ref{chap:requirements_spec}), where it is stated that Snowman shall exhibit a
sublinear increase in processing time in correlation to the number of rules in
the update (where the number of rules depicts the update size). Results are ranked
from strong correlation to no correlation, where stronger is better.
% TODO: remeber to explain, either here or in the rec. spec., why this is an
% important measure, e.g. no coorelation might imply programming errors.
\paragraph{Test cases:}
This test uses the same test cases as Test 1-1.

\subsubsection{Test 1-3}
\paragraph{Problem: Does raising the ratio between files and rules increase
processing time?}
Updates can come with the rules either gathered in one file, or distributed over 
several files. This question is asked to examine the difference in processing
time for these two scenarios. Results are ranked from no difference to great
difference, where less difference is better.
\paragraph{Test cases:}
This test compares the \textit{SYN-1000} test case from Test 1 with a
test case with the same number of rules in \textit{one} file
(\textit{SYN-1000-B}).

\subsubsection{Test 1-4}
\paragraph{Problem: Is the processing time shorter when the update contains no
new rules?}
This is a test for update processing when the update is not new (i.e. it has
been completely processed by the system before). To minimize system load, it is
essential that already processed updates are (in practice, they should be)
neglected by the system.
Results are ranked from shortest to longest processing time, where shorter is better.
\paragraph{Test cases:}
This is a compound test where a test case is run two times in succession. The
\textit{SYN-1000} test case from Test 1 is run first to populate the database,
before it is run again (\textit{SYN-1000-C}).

\subsubsection{Test 1-5}
\paragraph{Problem: Are the processing times for two updates identical when
only revision numbers are changed between the updates?}
This question measures how one small change in all rules affects procesing time.
This test runs the \textit{SYN-1000} update to populate the database, and then
the \textit{SYN-1000-D}, which contains a revision number change for all rules.
Processing time is expected to be faster for the second update, since the
database will already be populated with the majority of the data. Results from
this test will be ranked by the difference in processing time between the first
and the second updates, where greater difference is better.

\subsubsection{Test 1-6}
\paragraph{Problem: Are the processing times for two updates identical when only
one rule has changed between the updates?} This question seeks to determine how
the software handles an update when there is only one change from the previous
update. This is an important factor to consider, since an update usually
contains only a small amount of changes from previous updates. Results are
ranked from short processing time for second update to no difference in
processing time, where shorter time is better.
\paragraph{Test cases:}
This test suite contains two test cases with an update from a real source. The
\textit{SF28061} is run first to populate the database before \textit{SF28061-B}
is run, which contains one change (one of the rules has an updated revision
number).

%TODO: when analyzing results, be sure to highligt that a completely new
% update is run once, and subsequent updates often contain a small amount 
% of changes. 

\subsection{Rule distribution tests}
These tests use exactly the same test suites as update processing tests, the
only difference being that an update from a source is replaced by
synchronisation to a client.

\subsection{Interface loading tests}
These tests will evaluate the speed of the user interface by listing rules and
rule sets in a web browser. Tests are timed by using the development tools of
the Chrome browser, which contains a simple timer. Both tests use the SYN-10000
package.

\subsubsection{Test 3-1}
\paragraph{Problem: How quickly does the user interface list rules?}
\paragraph{Test cases:}
This test uses one test case, called {\em list rules}, where 100 rules are
listed in the user interface.

\subsubsection{Test 3-2}
\paragraph{Problem: How quickly does the user interface list rule sets?}
\paragraph{Test cases:}
This test uses one test case, called {\em list rule sets}, where all rule sets
are listed in the user interface.
